\documentclass[english,twoside, openright, hidelinks, 11 pt, class=report,crop=false]{standalone}
\input{../../preamb}
\input{../bm}
\setcounter{chapter}{1}
\setcounter{section}{1}
\begin{document}
\op{opgvekuperp}	
\st{
Gitt en vektor $\textbf{w}$. Vi definerer $\textbf{w}_p$ som vektoren som
\begin{itemize}
	\item står vinkelrett på $\textbf{w}$.
	\item peker til venstre med hensyn til retningen til \textbf{w}.
	\item har lik lengde som $\textbf{w}$.
\end{itemize}
Hvis ${\textbf{w}=[x, y]}$, kan det vises at ${\textbf{w}_p=[-y, x]}$, og at $-\textbf{w}_p$ peker til høyre for $\textbf{w}$.
\fig{uandup}
} \regv \regv 

Gitt vektorene $\textbf{u}=[a, b]$, $\textbf{u}_p=[-b, a]$, $\textbf{v}=[c, d]$ og $\textbf{v}_p=[-d, c]$. Vis at 
\abc{
\item $\textbf{u}\cdot\textbf{v}_p+\textbf{u}_p\cdot\textbf{v}=0$
\item $(\textbf{u}_p)_p=-\textbf{u}$
\item $(\textbf{u}+\textbf{v})_p=\textbf{u}_p+\textbf{v}_p$
}
\sv
\abc{
\item $[a, b]\cdot[-d, c] + [-b, a]\cdot[c, d] =-ad+bc-bd+ad=0$
\item $(\textbf{u}_p)_p=([-b, a])_p=[-a, -b]=-\textbf{u}$.
\item \phantom{} \vs \vs
\begin{align*}
	(\textbf{u}+\textbf{v})_p&=([a+c, b+d])_p\\
	&=[-(b+d), a+c]\\
	&=[-b, a]+[-d, c]\\ 
	&=\textbf{u}_p+\textbf{v}_p
\end{align*}
}
\newpage
\op{opgvekbotema}
Den røde og den grønne firkanten er begge et kvadrat. Vi setter\os $\textbf{v}=\vv{AC}$ og $\textbf{u}=\vv{CB}$. $F$ er midtpunktet på $DE$.\os
\abc{
\item Finn koordinatene til $F$ uttrykt ved $\textbf{u}$, $\textbf{v}$, $\textbf{u}_p$ og $\textbf{v}_p$.
\item Beskriv den eventuelle forflytningen til $F$ hvis $C$ forskyves i en vilkårlig retning. (Det holder å se på tilfellet hvor $C$ ligger over $AB$).
}
\fig{opgvekbottema}

\sv

\fig{opgvekbottema_lf}
Vi innfører et koordinatsystem med origo i $A$. Da er $D=\textbf{v}_p$. Videre er $E=\textbf{u}+\textbf{v}+\textbf{u}_p$. Videre er $\vv{DE}=-\textbf{v}_p+\textbf{v}+\textbf{u}+\textbf{u}_p$ Dermed er
\[ F = D+\frac{1}{2}\vv{DE}=\frac{1}{2}(\textbf{v}_p+\textbf{v}+\textbf{u}+\textbf{u}_p)=\frac{1}{2}(\textbf{u}+\textbf{v}+\textbf{u}_p+\textbf{v}_p) \]
$\textbf{u}+\textbf{v}$ og $\textbf{u}_p+\textbf{v}_p$ er ortogonale og har lik lengde. Så lenge $\textbf{u}+\textbf{v}$ ikke endres, vil altså ikke $F$ forflyttes. En forskyvning av $C$ medfører ingen endring av $\textbf{u}+\textbf{v}$, og dermed er poisjonen til $F$ uavhengig av posisjonen til $C$.

\newpage
\op{opgveknapl}
De blå linjestykkene inngår i likesidede trekanter på utsiden av $\triangle ABC$. $D$, $E$ og $F$ er de respektive midtpunktene til de likesidede trekantene. \os

Vis at $\triangle DEF$ er likesidet.
\fig{opgveknapl}.

\sv
Vi setter $\vv{AC}=2\textbf{v}$ og $\vv{AB}=2\textbf{u}$. Da er $\vv{BC}=2\textbf{v}-2\textbf{u}$. Videre har vi at
\begin{align*}
	D&=\textbf{v}+k\textbf{v}_p \vn
	E&=\textbf{u}-k\textbf{u}_p \vn
	F&=\textbf{u}+\textbf{v}-k(\textbf{v}_p-\textbf{u}_p)
\end{align*} 
hvor $k=\frac{1}{\sqrt{3}}$. Altså er
\begin{align*}
	\vv{DF}&=\textbf{u}+k(\textbf{u}_p-2\textbf{v}_p) \\
	\vv{DE}&=\textbf{u}-\textbf{v}-k(\textbf{u}_p+\textbf{v}_p) \\
	\vv{EF}&=\textbf{v}+k(2\textbf{u}_p-\textbf{v}_p)
\end{align*} 
Da er
\begin{align*}
DF^2&=\textbf{u}^2+2k\textbf{u}(\textbf{u}_p-2\textbf{v}_p)+k^2(\textbf{u}_p-2\textbf{v}_p)^2\\
&=\textbf{u}^2-4k\textbf{u}\cdot\textbf{v}_p+k^2(\textbf{u}_p^2-4\textbf{u}_p\cdot\textbf{v}_p+4\textbf{v}_p^2) \\
&=\frac{4}{3}(\textbf{u}^2+\textbf{v}^2-\textbf{u}_p\cdot\textbf{v}_p)-4k\textbf{u}\cdot\textbf{v}_p
\end{align*}
\begin{align*}	DE^2&=(\textbf{u}-\textbf{v})^2-2k(\textbf{u}-\textbf{v})(\textbf{u}_p+\textbf{v}_p)+k^2(\textbf{u}_p+\textbf{v}_p)^2 \\
&=\textbf{u}^2-2\textbf{u}\cdot\textbf{v}+\textbf{v}^2-2k(\textbf{u}\cdot\textbf{v}_p-\textbf{u}_p\cdot\textbf{v})+k^2(\textbf{u}_p^2+2\textbf{u}_p\cdot\textbf{v}_p+\textbf{v}_p^2)\\
&=\frac{4}{3}(\textbf{u}^2+\textbf{v}^2-\textbf{u}_p\cdot\textbf{v}_p)-4k\textbf{u}\cdot\textbf{v}_p
\end{align*}

\begin{align*}
	EF^2 &= \textbf{v}^2+2k\textbf{v}(2\textbf{u}_p-\textbf{v}_p)+k^2(2\textbf{u}_p-\textbf{v}_p)^2 \\
	&=\textbf{v}^2+4k\textbf{u}_p\cdot\textbf{v}+k^2(4\textbf{u}_p-4\textbf{u}_p\textbf{v}_p+\textbf{v}_p^2) \\
	&=\frac{4}{3}(\textbf{u}^2+\textbf{v}^2-\textbf{u}_p\cdot\textbf{v}_p)-4k\textbf{u}\cdot\textbf{v}_p
\end{align*}
\newpage
\op{opgvekvanbl}
De fargede firkantene er kvadrater. $E$, $F$, $G$ og $H$ er de respektive midtpunktene til kvadratene. Vis at $EG=FH$ og $\vv{EG}\perp\vv{FH}$.
\fig{opgvekvanbl}

\textbf{Svar}
\fig{opgvekvanbl_lf}
Vi innfører et koordinatsystem med origo i $A$, og setter $\vv{DA}=2\textbf{v}$, $\vv{AB}=2\textbf{u}$, og $\vv{BC}=2\textbf{w}$ (i figuren er $\textbf{s}=\textbf{u}+\textbf{v}+\textbf{w}$). Da er
\begin{align*}
	E&=-\textbf{v}-\textbf{v}_p \\
	F&=\textbf{u}-\textbf{u}_p \\
	G &= 2\textbf{u}+\textbf{w}-\textbf{w}_p\\
	H&=\textbf{u}+\textbf{w}-\textbf{v}+\textbf{u}_p+\textbf{v}_p+\textbf{w}_p
\end{align*}
Dermed er 
\begin{align*}
\vv{EG}&=2\textbf{u}+\textbf{w}-\textbf{w}_p+\textbf{v}+\textbf{v}_p=2\textbf{u}+\textbf{w}+\textbf{v}+\textbf{v}_p-\textbf{w}_p \\
\vv{FH}&=\textbf{w}-\textbf{v}+2\textbf{u}_p+\textbf{v}_p+\textbf{w}_p
\end{align*}
Vi har at
\[ \vv{FH}_p=\textbf{w}_p-\textbf{v}_p-2\textbf{u}-\textbf{v}-\textbf{w}=-\vv{EG} \]
Dette betyr at $\vv{FH}$ og $\vv{EG}$ er ortogonale. Videre er
\[ FH=|\vv{FH}_p|=|-\vv{EG}|=EG \]

\newpage
\op{opgvekuperp}
\st{
	Given a vector $\textbf{w}$. We define $\textbf{w}_p$ as the vector that
	\begin{itemize}
		\item is perpendicular to $\textbf{w}$.
		\item points to the left with respect to the direction of $\textbf{w}$.
		\item has the same length as $\textbf{w}$.
	\end{itemize}
	If ${\textbf{w}=[x, y]}$, it can be shown that ${\textbf{w}_p=[-y, x]}$, and that $-\textbf{w}_p$ points to the right of $\textbf{w}$.
	\fig{uandup}
} \regv \regv

Given the vectors $\textbf{u}=[a, b]$, $\textbf{u}_p=[-b, a]$, $\textbf{v}=[c, d]$, and $\textbf{v}_p=[-d, c]$. Show that
\abc{
	\item $\textbf{u}\cdot\textbf{v}_p+\textbf{u}_p\cdot\textbf{v}=0$
	\item $(\textbf{u}_p)_p=-\textbf{u}$
	\item $(\textbf{u}+\textbf{v})_p=\textbf{u}_p+\textbf{v}_p$
}
\textbf{Answer} \\
\abc{
	\item $[a, b]\cdot[-d, c] + [-b, a]\cdot[c, d] =-ad+bc-bd+ad=0$
	\item $(\textbf{u}_p)_p=([-b, a])_p=[-a, -b]=-\textbf{u}$.
	\item \phantom{} \vs \vs
	\begin{align*}
		(\textbf{u}+\textbf{v})_p&=([a+c, b+d])_p\\
		&=[-(b+d), a+c]\\
		&=[-b, a]+[-d, c]\\
		&=\textbf{u}_p+\textbf{v}_p
	\end{align*}
}

\newpage
\subsection*{Problem: Bottema's theorem}

The red and the green quadrilateral are both squares. $F$ is the midpoint of $DE$. Let $\textbf{v}=\vv{AC}$ and $\textbf{u}=\vv{CB}$. Given a vector $\textbf{w}$, we define $\textbf{w}_p$ as the vector perpendicular to $\textbf{w}$, pointing to the left with respect to the direction of $\textbf{w}$, and with equal length as $\textbf{w}$.\os
\begin{enumerate}[label=\alph*)]
	\item Find the coordinates of $F$ expressed in terms of $\textbf{u}$, $\textbf{v}$, $\textbf{u}_p$, and $\textbf{v}_p$.
	\item Describe the possible displacement of $F$ if $C$ is shifted in an arbitrary direction. (It is sufficient to consider the case where $C$ is above $AB$).
\end{enumerate}
\fig{opgvekbottema}

\subsection*{Solution}

\fig{opgvekbottema_lf}

We introduce a coordinate system with the origin at $A$. Then $D=\textbf{v}_p$. Furthermore, $E=\textbf{u}+\textbf{v}+\textbf{u}_p$. Furthermore, $\vv{DE}=-\textbf{v}_p+\textbf{v}+\textbf{u}+\textbf{u}_p$. Thus,
\[ F = D+\frac{1}{2}\vv{DE}=\frac{1}{2}(\textbf{v}_p+\textbf{v}+\textbf{u}+\textbf{u}_p)=\frac{1}{2}(\textbf{u}+\textbf{v}+\textbf{u}_p+\textbf{v}_p) \]

$\textbf{u}+\textbf{v}$ and $\textbf{u}_p+\textbf{v}_p$ are orthogonal and have the same length. As long as $\textbf{u}+\textbf{v}$ does not change, $F$ will not be displaced. A displacement of $C$ does not change $\textbf{u}+\textbf{v}$, and thus the position of $F$ is independent of the position of $C$.


\newpage
\textbf{Problem: van Aubel's theorem} \\
The colored quadrilaterals are squares. $E$, $F$, $G$, and $H$ are the respective midpoints of the squares. Show that $EG = FH$ and $\vv{EG} \perp \vv{FH}$.
\fig{opgvekvanbl}

\textbf{Answer}
\fig{opgvekvanbl_lf}
We introduce a coordinate system with the origin at $A$, and set $\vv{DA} = 2\textbf{v}$, $\vv{AB} = 2\textbf{u}$, and $\vv{BC} = 2\textbf{w}$ (in the figure, $\textbf{s} = \textbf{u} + \textbf{v} + \textbf{w}$). Then,
\begin{align*}
	E &= -\textbf{v} - \textbf{v}_p \\
	F &= \textbf{u} - \textbf{u}_p \\
	G &= 2\textbf{u} + \textbf{w} - \textbf{w}_p \\
	H &= \textbf{u} + \textbf{w} - \textbf{v} + \textbf{u}_p + \textbf{v}_p + \textbf{w}_p
\end{align*}
Thus,
\begin{align*}
	\vv{EG} &= 2\textbf{u} + \textbf{w} - \textbf{w}_p + \textbf{v} + \textbf{v}_p = 2\textbf{u} + \textbf{w} + \textbf{v} + \textbf{v}_p - \textbf{w}_p \\
	\vv{FH} &= \textbf{w} - \textbf{v} + 2\textbf{u}_p + \textbf{v}_p + \textbf{w}_p
\end{align*}
We have
\[
\vv{FH}_p = \textbf{w}_p - \textbf{v}_p - 2\textbf{u} - \textbf{v} - \textbf{w} = -\vv{EG}
\]
This means that $\vv{FH}$ and $\vv{EG}$ are orthogonal. Furthermore,
\[ FH = |\vv{FH}_p| = |-\vv{EG}| = EG \]

\newpage
$A$, $B$, $C$ og $D$ ligger på sirkelen, og $\angle ACB=\angle BAC=v$.\os

Vis at 
\[ BD = \frac{CD+DA}{2\cos v} \] 
\fig{opggeoisoabc}

\sv
$\angle BDC$ og $\angle ADB$ spenner over samme bue som henholdsvis $\angle ACB$ og $\angle BAC$. Alle disse er periferivinkler, og da $\angle ACB=\angle BAC$, har de alle like størrelse. Videre følger det at $BC=AB$. Av cosinussetingen har vi da at
\begin{align*}
	CD^2+BD^2-2CD\cdot BD\cos v &= DA^2+BD^2-2DA\cdot BD\cos v \\
	CD^2-DA^2 &= 2BD\cos v(CD-DA) \\
	(CD+DA)(CD-DA)&= 2BD\cos v(CD-DA)
\end{align*}
Antatt at $CD-DA\neq0$, er 
\[ BD=\frac{CD+DA}{2\cos v} \]
I tilfellet hvor ${CD=DA}$, er $BD$ en diameter i sirkelen. Da har vi av Thales' setning omvendt at $\angle DCB=90^\circ$. Altså er
\begin{align*}
	BC &= \frac{CD}{\cos v}  = \frac{CD+DA}{2\cos v} \br
\end{align*}
\newpage
$A$, $B$, $C$ and $D$ lies on the circle, and
$\angle ACB=\angle BAC=v$.\os 

Show that
\[ BD = \frac{CD+DA}{2\cos v} \] 
\fig{opggeoisoabc}
\newpage
$\triangle ABC$ er likesidet, og $\triangle ACD\cong\triangle BAE$.
\abc{
\item Avgjør hvilke av disse vinklene som har lik verdi. Bestem vinkel-verdien om det lar seg gjøre.
\[\angle BDC, \angle ADB, \angle CFD, \angle ACE, \angle BEC, \angle EFB\text{ og } \angle CAD \]
\item Vis at $BD=CD+DA$.
}
\fig{opggeoisab}
\sv
\abc{
\item $\angle BDC=\angle BEC=\angle BAC=60^\circ$ fordi de er periferivinkler som spenner over samme bue. Tlsvarende kan vi finne at $\angle ADB=60^\circ$. $\angle CAD=\angle ACE$ fordi de er periferivinkler som spenner over sirkelsegmenter med lik lengde. Av vinkelsummen til $\angle ADC$ er $\angle DCA=60-\angle CAD$. Dermed er
\[ \angle DCF=\angle DCA+ACE=60^\circ \]
Altså er $\triangle CDF$ likesidet. $\angle EFB=\angle CFD$ fordi de er toppvinkler, og dermed er også $\triangle FEB$ likesidet. Følgelig har vi at
\[ \angle BDC=\angle ADB=\angle CFD=\angle BEC=\angle EFB=60^\circ \]
og 
\[ \angle CAD=\angle ACE \]
\item Da $\triangle CDF$ og $\triangle FEB$ er likesidet, har vi at
\[ BD=BF+FD=BE+CD=AD+CD \]
}

\end{document}
